\usepackage[utf8]{inputenc}
\usepackage[T1]{fontenc}					% Proper font encoding for Indonesian/Latin characters
\usepackage[bahasa,english]{babel}			% Bilingual: Indonesian (bahasa) + English
\usepackage{verbatim}						% To use \begin{comment} text to comment \end{comment} or inserting codes \begin{verbatim}
%\usepackage{minted}						    % Insert code programming inline
\usepackage{listings}
\usepackage{xcolor}
\usepackage[titles]{tocloft}
\usepackage{cite}							% To cite your references
\usepackage{graphicx}						% Provide features to customize graphics
\usepackage{draftwatermark}					% To add a watermark : default = DRAFT
\usepackage{setspace}						% Provide support for setting the spacing between lines
\usepackage{titlesec}						% To modify the seciton titles
\usepackage{nomencl}						% To create nomenclature list
\makenomenclature							% Activate the nomenclature
\usepackage{acronym}						% To create Acronyms list
\usepackage{booktabs}
\usepackage{footnote}
\usepackage{caption}						% To manipulate captions for Figures and Tables
\DeclareCaptionFont{white}{\color{white}}
\DeclareCaptionFormat{listing}{\colorbox{gray}{\parbox{\textwidth}{#3}}}
\captionsetup[lstlisting]{format=listing,labelfont={bf,white},textfont={bf,white}}
\captionsetup{font=large}
\usepackage{subfig}
\usepackage{floatrow}						% To add a note under captions of Tables and Figures
\floatsetup[table]{capposition=top,font=large}			% To force the Tables captions back on top
\usepackage[flushleft]{threeparttable}		% To create tables with three sections (title, table, note)
\usepackage{array}							% To handle tabular and array environments
\usepackage{chngcntr}	
\usepackage{multirow}	
\usepackage{amsmath}

% To manipulate counters for Figures and Tables
\usepackage{pdfpages}						% To Include PDFS
\counterwithout{equation}{chapter} 			% Remove the chapter number in equations numbering
% \counterwithin{equation}{chapter}  		% Add the chapter number if you prefer but I think it's not correct with given format
\counterwithin{table}{chapter} 			% Remove the chapter number in tables numbering
% set the default margins of the layout, it will be overwritten at the beginning for first pages
\usepackage[a4paper,top=25mm,bottom=25mm,left=20mm,right=20mm,bindingoffset=10mm]{geometry}
\pagestyle{plain}
% Redefinition of the chapter command
\makeatletter
\def\@chapter[#1]#2{\ifnum \c@secnumdepth >\m@ne
                       \if@mainmatter
                         \refstepcounter{chapter}%
                         \typeout{\@chapapp \space \thechapter.}%
                         \addcontentsline{toc}{chapter}%
                                   {\protect\numberline{\thechapter}#1}%
                       \else				
                         \addcontentsline{toc}{chapter}{#1}%
                       \fi
                    \else                      
                      \addcontentsline{toc}{chapter}{#1}%
                    \fi
                    \chaptermark{#1}%
                    \if@twocolumn
                      \@topnewpage[\@makechapterhead{#2}]%
                    \else
                      \@makechapterhead{#2}%
                      \@afterheading
                    \fi}
\makeatother

\titleformat{\chapter}[display]
{\fontsize{30}{10}\selectfont 
\normalfont\bfseries\centering}
{\fontsize{30}{10}\selectfont \chaptertitlename\ \thechapter}{0.5em}{}

\makeatletter \@removefromreset{figure}{chapter}\makeatother

%% Bilingual section headings (Indonesian primary)
\renewcommand{\contentsname}{Daftar Isi}
\renewcommand{\listfigurename}{Daftar Gambar}
\renewcommand{\listtablename}{Daftar Tabel}
\renewcommand{\bibname}{Daftar Pustaka}
\renewcommand{\appendixname}{Lampiran}
\renewcommand{\thechapter}{\Roman{chapter}}
\renewcommand{\thesection}{\arabic{chapter}.\arabic{section}}

\titleformat{\section}{\normalfont\Large\bfseries}{\thesection}{1em}{\setstretch{0.1}}

\renewcommand{\emph}[1]{``#1''}
\renewcommand{\thefigure}{\arabic{chapter}.\arabic{figure}}
\renewcommand{\thetable}{\arabic{chapter}.\arabic{table}}

\newcommand\T{\rule{0pt}{2.6ex}}       % Top strut - used in Tables to add space on Top of the cell
\newcommand\B{\rule[-1.2ex]{0pt}{0pt}} % Bottom strut - used in Tables to add space on Bottom of the cell

\renewcommand*\cftchapnumwidth{3em}
\renewcommand{\cftchapdotsep}{\cftdotsep}% Chapters should use dots in ToC

%%%% Language switching short commands %%%%%%%%%%%%%%%%%%%%%%%
\newcommand{\IDon}{\begin{otherlanguage}{bahasa}} %% Opens Indonesian writing
\newcommand{\IDoff}{\end{otherlanguage}}           %% Closes Indonesian writing
\newcommand{\ENon}{\begin{otherlanguage}{english}} %% Opens English writing
\newcommand{\ENoff}{\end{otherlanguage}}           %% Closes English writing
